% ── CAPÍTULO 3: INSTALACIÓN ──────────────────────────────────
\section{Instalación Paso a Paso}
\label{sec:instalacion}

\subsection{Método 1: Con pip (estándar)}

\subsubsection{Paso 1 — Clonar el repositorio}

\begin{lstlisting}[style=shell, caption={Clonar el repositorio desde GitHub.}]
git clone https://github.com/SebastianG210/proyectoAlgoritmos.git
cd proyectoAlgoritmos
\end{lstlisting}

\subsubsection{Paso 2 — Crear un entorno virtual}

\begin{lstlisting}[style=shell, caption={Crear y activar el entorno virtual.}]
# Crear el entorno
python -m venv .venv

# Activar en Windows PowerShell
.venv\Scripts\Activate.ps1

# Activar en Linux/macOS
source .venv/bin/activate
\end{lstlisting}

\begin{aviso}
En Windows, si aparece el error \textit{«la ejecución de scripts está deshabilitada»}, ejecute primero:\\
\lstinline|Set-ExecutionPolicy -ExecutionPolicy RemoteSigned -Scope CurrentUser|
\end{aviso}

\subsubsection{Paso 3 — Instalar dependencias}

\begin{lstlisting}[style=shell, caption={Instalar el proyecto en modo editable.}]
pip install -e .
\end{lstlisting}

\subsubsection{Paso 4 — Verificar la instalación}

\begin{lstlisting}[style=shell, caption={Verificar que las dependencias están instaladas correctamente.}]
python -c "import numpy, plotly, pandas; print('OK')"
\end{lstlisting}

Si el resultado es \lstinline|OK|, la instalación fue exitosa.

\subsection{Método 2: Con uv (recomendado)}

\begin{lstlisting}[style=shell, caption={Instalación rápida con uv.}]
# Instalar uv (si no está disponible)
pip install uv

# Clonar e instalar en un solo paso
git clone https://github.com/SebastianG210/proyectoAlgoritmos.git
cd proyectoAlgoritmos
uv sync
\end{lstlisting}

\begin{tip}
\lstinline|uv sync| lee el archivo \lstinline|pyproject.toml| y crea automáticamente el entorno virtual con todas las dependencias, sin pasos adicionales.
\end{tip}

\subsection{Estructura de Archivos tras la Instalación}

\begin{lstlisting}[caption={Estructura relevante del proyecto después de instalar.}]
proyectoAlgoritmos/
+-- src/
|   +-- shared/
|   |   +-- input/       <- Archivos CSV de redes de prueba
|   |   \-- output/      <- Resultados generados aqui
|   \-- strategies/      <- Algoritmos KGeoMIP y KQNodes
+-- pyproject.toml       <- Configuracion de dependencias
\-- document/            <- Este manual
\end{lstlisting}

\subsection{Actualización del Software}

\begin{lstlisting}[style=shell, caption={Actualizar a la versión más reciente.}]
git pull origin master
# Con pip:
pip install -e . --upgrade
# Con uv:
uv sync
\end{lstlisting}
